\scanpage{113}
\mspage{28}
% archive sheet number: 112
\deleted{une suite exacte canonique}

\[
\deleted{0\longrightarrow\bigwedge^i\Lambda^{r-1}\otimes}
\]

\[
\begin{aligned}
H^1(U,\Lambda)&\simeq\coprod_{1\leq i\leq r}H^1(U_i,\Lambda)\simeq\coprod_{1\leq i\leq r-1}H^1(U_i,\Lambda)\times H^1(U_r,\Lambda)\\
&\simeq H^1(U',\Lambda)\times\mu_n(X)^{-1}
\end{aligned}
\]

donc une \deleted{suite exacte} \inserted{décomposition}

\[
\bigwedge^iH^1(U,\Lambda)\simeq\bigwedge^iH^1(U',\Lambda)\times\left(\bigwedge^{i-1}H^1(U',\Lambda)\otimes\mu_n(X)^{-1}\right)
\]

\[
\deleted{0\longrightarrow\bigwedge^iH^1(U',\Lambda)}
\tag*{\deleted{(4.5.7)}}
\]

\deleted{\unsure{qui fournit}} qui définit une suite exacte

\[
0\longrightarrow\bigwedge^iH^1(U',\Lambda)\longrightarrow\bigwedge^iH^1(U,\Lambda)\longrightarrow\left(\bigwedge^{i-1}H^1(U',\Lambda)\right)(-1)\longrightarrow0
\tag*{(4.5.7)}
\]

On \deleted{\unsure{obtient}} \unsure{construit} \inserted{par cup-produit} des flèches \deleted{\unsure{dans|can.}} des termes de (4.5.7) dans les termes correspondants de (4.5.6). Nous admettrons qu'on obtient ainsi un hom. de suites exactes (c'est trivial pour le premier carré, et semble demander une vérification pour le deuxième). L'\deleted{\unsure{autre|actuelle}} hypothèse de récurrence (appliquée à \deleted{$\unsure{U'=X-\bigcup D_i}$ \unsure{resp. à} $D'_r$} $U'$ resp. à $D'_r$) implique que les flèches verticales extrêmes sont \deleted{\unsure{des isom}} \inserted{bijectives}, donc aussi la flèche verticale médiane, ce qui achève la démonstration.
